% Présentation de ThesisManager (mod_these et mod_thesesynthesis).
% Compilation : pdflatex presentation-these.tex (deux fois), depuis docs/presentation.
\documentclass[aspectratio=169,11pt]{beamer}

\usepackage[T1]{fontenc}
\usepackage[utf8]{inputenc}
\usepackage{lmodern}
\usepackage[french]{babel}
\usepackage{microtype}
\usepackage{booktabs}
\usepackage{tabularx}
\usepackage{tikz}
\usetikzlibrary{arrows.meta,positioning,shapes.geometric}

\graphicspath{{../guides/images/}}

% Couleurs : celles du guide (bleu Moodle).
\definecolor{moodle}{HTML}{0F6CBF}
\definecolor{orange}{HTML}{F98012}
\definecolor{gris}{HTML}{5F6B76}
\definecolor{fond}{HTML}{F3F6F9}

\usetheme{default}
\setbeamercolor{structure}{fg=moodle}
\setbeamercolor{frametitle}{fg=moodle}
\setbeamercolor{title}{fg=moodle}
\setbeamercolor{alerted text}{fg=orange}
\setbeamerfont{frametitle}{series=\bfseries,size=\Large}
\setbeamertemplate{navigation symbols}{}
\setbeamertemplate{itemize items}{\color{moodle}$\blacktriangleright$}
\setbeamertemplate{footline}{%
  \hfill{\color{gris}\scriptsize ThesisManager \quad \insertframenumber/\inserttotalframenumber}\hspace*{1em}\vspace*{0.6em}}

% Capture bordée, tronquée en bas (trim en pixels de l'image, à l'échelle 1px = 1bp).
\newcommand{\capture}[3][0]{%
  {\setlength{\fboxsep}{0pt}\color{gris!50}\fbox{\includegraphics[#2,trim=0 #1 0 0,clip]{#3}}}}

% Étape du circuit.
\tikzset{
  etape/.style={rectangle,rounded corners=3pt,draw=moodle,fill=fond,thick,align=center,
    font=\footnotesize,minimum height=1.1cm,text width=2.25cm},
  acteur/.style={font=\scriptsize\bfseries,text=orange},
  fleche/.style={-{Stealth[length=2.5mm]},thick,moodle},
  retour/.style={-{Stealth[length=2mm]},gris,dashed},
}

\title{ThesisManager}
\subtitle{La gestion des thèses d'exercice vétérinaire dans Moodle,\\du dépôt du sujet au procès-verbal}
\author{Sébastien Lefebvre}
\date{Version 0.19.0 -- octobre 2026}

\begin{document}

\begin{frame}[plain]
  \titlepage
\end{frame}

% ---------------------------------------------------------------------------
\begin{frame}{Pourquoi un module dédié ?}
  \begin{columns}[T]
    \begin{column}{0.48\textwidth}
      \textbf{Aujourd'hui}
      \begin{itemize}
        \item Des formulaires PDF et Word qui circulent par courriel.
        \item Des signatures manuscrites à recueillir une par une.
        \item Un suivi dispersé entre l'étudiant, le directeur, la DEVE et la bibliothèque.
        \item Des relances faites à la main.
      \end{itemize}
    \end{column}
    \begin{column}{0.48\textwidth}
      \textbf{Avec ThesisManager}
      \begin{itemize}
        \item Un \alert{dossier unique} par étudiant, de la fiche de sujet au procès-verbal.
        \item Chacun voit ce qui l'attend et agit au bon moment.
        \item Les documents officiels sont \alert{générés et signés} dans Moodle.
        \item Courriels, relances et historique sont automatiques.
      \end{itemize}
    \end{column}
  \end{columns}
  \vfill
  \begin{center}
    \small\color{gris}Il suit le document de travail « Circuit de dépôt du sujet et de la soutenance ».
  \end{center}
\end{frame}

% ---------------------------------------------------------------------------
\begin{frame}{Deux activités Moodle complémentaires}
  \begin{columns}[T]
    \begin{column}{0.48\textwidth}
      \begin{block}{Gestion des thèses \texttt{(mod\_these)}}
        Une activité par promotion.
        \begin{itemize}
          \item Dossier de l'étudiant : sujet, soutenance, jury.
          \item Validations, renvois, signatures.
          \item Documents de soutenance, procès-verbal, registre.
        \end{itemize}
      \end{block}
    \end{column}
    \begin{column}{0.48\textwidth}
      \begin{block}{Suivi des thèses \texttt{(mod\_thesesynthesis)}}
        Une seule activité pour toutes les promotions.
        \begin{itemize}
          \item Vue unique des dossiers pour directeurs, DEVE, bibliothèque et jurys.
          \item Paramètres du circuit, modèles des documents, grille et questionnaire communs.
          \item Annuaire du jury, recherche, planning.
        \end{itemize}
      \end{block}
    \end{column}
  \end{columns}
\end{frame}

% ---------------------------------------------------------------------------
\begin{frame}{Les acteurs, sans rôle à créer}
  \small
  \begin{tabularx}{\textwidth}{@{}llX@{}}
    \toprule
    \textbf{Acteur} & \textbf{Rôle Moodle} & \textbf{Ce qu'il fait} \\
    \midrule
    Étudiant & Étudiant & propose son sujet, dépose sa soutenance \\
    Directeur de thèse & Enseignant non éditeur & valide, rédige son rapport, signe \\
    DEVE & Enseignant & enregistre, valide, génère les documents \\
    Agents DEVE signataires & Enseignant, désignés & signent pour la DEVE et, par délégation, pour la DG \\
    Bibliothèque & inscrite, désignée agent & notice du sujet, conformité du manuscrit \\
    Jury, signataires extérieurs & sans compte & accès et signature par lien personnel \\
    \bottomrule
  \end{tabularx}
  \vfill
  \small Les agents se \alert{choisissent dans une liste} des personnes inscrites : aucun rôle spécifique, aucune saisie d'identifiant.
\end{frame}

% ---------------------------------------------------------------------------
\begin{frame}{Phase 1 : le sujet}
  \centering
  \begin{tikzpicture}[node distance=0.55cm,every node/.append style={}]
    \tikzset{etape/.append style={text width=2.0cm}}
    \node[etape] (e1) {Fiche de sujet\\signée};
    \node[etape,right=of e1] (e2) {Validation\\du sujet};
    \node[etape,right=of e2] (e3) {Enregistrement\\du sujet};
    \node[etape,right=of e3] (e4) {Notice\\documentaire};
    \node[etape,right=of e4,fill=moodle!15] (e5) {Dossier\\actif};
    \node[acteur,above=0.15cm of e1] {Étudiant};
    \node[acteur,above=0.15cm of e2] {Directeur};
    \node[acteur,above=0.15cm of e3] {DEVE};
    \node[acteur,above=0.15cm of e4] {Bibliothèque};
    \draw[fleche] (e1) -- (e2);
    \draw[fleche] (e2) -- (e3);
    \draw[fleche] (e3) -- (e4);
    \draw[fleche] (e4) -- (e5);
    \draw[retour] (e3.south) to[bend left=25] node[below,font=\scriptsize] {renvoi motivé} (e1.south);
  \end{tikzpicture}
  \vfill
  \begin{itemize}
    \item Chaque étape peut \alert{renvoyer} le dossier avec un motif ; à sa nouvelle soumission, il revient directement à l'étape qui l'a renvoyé.
    \item Le dossier reste ensuite actif, parfois plusieurs années ; un \alert{changement de sujet} archive l'ancien et relance la phase 1.
  \end{itemize}
\end{frame}

% ---------------------------------------------------------------------------
\begin{frame}{Phase 2 : la soutenance}
  \centering
  \begin{tikzpicture}[node distance=0.45cm]
    \node[etape] (s1) {Questionnaire,\\manuscrit,\\date, jury};
    \node[etape,right=of s1] (s2) {Rapport\\et grille,\\validation};
    \node[etape,right=of s2] (s3) {Dossier complet,\\n° DEVE,\\jury vérifié};
    \node[etape,right=of s3] (s4) {Conformité\\du manuscrit};
    \node[etape,right=of s4] (s5) {Convocations,\\documents,\\transmission};
    \node[etape,below=1.0cm of s5,fill=moodle!15] (s6) {Procès-verbal,\\dossier clos};
    \node[acteur,above=0.15cm of s1] {Étudiant};
    \node[acteur,above=0.15cm of s2] {Directeur};
    \node[acteur,above=0.15cm of s3] {DEVE};
    \node[acteur,above=0.15cm of s4] {Bibliothèque};
    \node[acteur,above=0.15cm of s5] {DEVE};
    \draw[fleche] (s1) -- (s2);
    \draw[fleche] (s2) -- (s3);
    \draw[fleche] (s3) -- (s4);
    \draw[fleche] (s4) -- (s5);
    \draw[fleche] (s5) -- (s6);
    \node[etape,left=1.2cm of s6,text width=4.3cm,draw=gris] (c) {\textbf{Corrections demandées} : manuscrit corrigé, validé par le directeur, la bibliothèque et la DEVE};
    \draw[retour] (s6) -- (c);
  \end{tikzpicture}
  \vfill
  \small Le directeur \alert{rédige son rapport dans un formulaire} (non communiqué à l'étudiant). Une thèse \alert{ajournée} repart pour un nouveau dépôt complet.
\end{frame}

% ---------------------------------------------------------------------------
\begin{frame}{L'étudiant suit son dossier}
  \begin{columns}[c]
    \begin{column}{0.66\textwidth}
      \capture{width=\linewidth}{01_etudiant_accueil}
      \vspace{0.4em}
      \begin{itemize}
        \item Une \alert{frise} indique l'étape atteinte et qui doit agir.
        \item Les boutons bleus sont les actions qui l'attendent.
        \item Brouillon possible avant chaque soumission.
        \item \alert{Annuaire des enseignants} en lecture seule : statut, matière, présidents de jury possibles et leur adresse.
      \end{itemize}
    \end{column}
    \begin{column}{0.28\textwidth}
      \centering
      \capture[300]{height=0.72\textheight}{104_frise_telephone}\\[0.3em]
      {\scriptsize\color{gris} Sur téléphone}
    \end{column}
  \end{columns}
\end{frame}

% ---------------------------------------------------------------------------
\begin{frame}{La DEVE pilote depuis un tableau de bord}
  \begin{columns}[T]
    \begin{column}{0.52\textwidth}
      \capture[640]{width=\linewidth}{18_tableau_de_bord_deve}
    \end{column}
    \begin{column}{0.44\textwidth}
      \begin{itemize}
        \item \alert{« À traiter »} : chaque dossier avec le bouton de l'action attendue.
        \item Actions groupées : relancer, générer les documents, prévenir le jury, enregistrer l'A6.
        \item Planning des soutenances (semaine, mois, PDF) et calendrier Moodle.
        \item Conflits de jury et de salle signalés.
        \item Récapitulatif hebdomadaire par courriel.
        \item Export du registre (CSV, Excel, ODS).
      \end{itemize}
    \end{column}
  \end{columns}
\end{frame}

% ---------------------------------------------------------------------------
\begin{frame}{Des documents officiels générés et signés}
  \begin{columns}[c]
    \begin{column}{0.55\textwidth}
      \capture{width=\linewidth}{42_documents_signes}
    \end{column}
    \begin{column}{0.41\textwidth}
      \begin{itemize}
        \item Fiche de sujet, agréments, fiche signalétique, bordereau ABES, rapport, grille, conclusions, contrat de diffusion, convocations, PV.
        \item Mise en forme de l'établissement, \alert{remplis automatiquement}.
        \item Chaque modèle est \alert{modifiable par la DEVE} (aperçu PDF, retour à l'original).
        \item Liste du corps enseignant en PDF, sur trois colonnes, avec la matière.
      \end{itemize}
    \end{column}
  \end{columns}
\end{frame}

% ---------------------------------------------------------------------------
\begin{frame}{Les signatures, sans papier}
  \begin{columns}[T]
    \begin{column}{0.47\textwidth}
      \capture[520]{width=\linewidth}{64_agent_mes_signatures}
    \end{column}
    \begin{column}{0.49\textwidth}
      \begin{itemize}
        \item Signature tracée, \alert{enregistrée} et reprise d'office.
        \item \alert{Mes signatures} : tous les documents en attente, toutes promotions.
        \item Signataires sans compte : lien personnel par courriel.
        \item Sur place, après vérification d'identité, ou par scan.
        \item La \alert{direction générale} est toujours signée par délégation, par les agents DEVE.
        \item Relances automatiques et bouton \og Relancer\fg{}.
      \end{itemize}
    \end{column}
  \end{columns}
\end{frame}

% ---------------------------------------------------------------------------
\begin{frame}{Intégrité et confidentialité}
  \begin{columns}[T]
    \begin{column}{0.48\textwidth}
      \textbf{Intégrité}
      \begin{itemize}
        \item Chaque signature garde l'\alert{empreinte} du contenu signé.
        \item Un document modifié après signature est invalidé.
        \item Document \alert{scellé} une fois toutes ses signatures réunies.
        \item Historique complet de chaque dossier.
      \end{itemize}
    \end{column}
    \begin{column}{0.48\textwidth}
      \textbf{Confidentialité}
      \begin{itemize}
        \item Thèse confidentielle : sujet masqué à qui n'a pas à le connaître.
        \item Embargo sur la diffusion.
        \item Liens à jeton : seule une empreinte est conservée, durée limitée.
        \item API vie privée de Moodle (export, effacement), sauvegarde et restauration.
      \end{itemize}
    \end{column}
  \end{columns}
\end{frame}

% ---------------------------------------------------------------------------
\begin{frame}{Jury et annuaire}
  \begin{itemize}
    \item \alert{Annuaire commun} à toutes les promotions : titre, statut, établissement, matière ; import CSV ou préremplissage depuis le cours.
    \item Le jury se compose depuis l'annuaire ; seuls les statuts admis peuvent présider (Pr, émérites, MCF HDR).
    \item Éligibilité des membres extérieurs vérifiée par la DEVE.
    \item Directeur de thèse membre \alert{non votant} ; au moins deux membres votants.
    \item Membres du jury : accès aux documents par lien ou depuis \og Mes jurys\fg{}.
    \item Étudiants : annuaire en lecture seule, filtre \og pouvant présider un jury\fg{}, adresse des présidents possibles.
  \end{itemize}
\end{frame}

% ---------------------------------------------------------------------------
\begin{frame}{Paramétrage par la DEVE, sans réglage du site}
  \begin{itemize}
    \item \alert{Paramètres du circuit} dans l'activité de suivi : établissement, signataires, agents, salles, durée des soutenances, contrat type.
    \item Encadré \og Réglages à compléter\fg{} tant qu'il manque quelque chose.
    \item \alert{13 courriels personnalisables}, avec variables.
    \item Modèles des documents, grille de conformité, questionnaire préalable au jury.
    \item \alert{Transfert des paramètres} d'une activité à une autre, ou par fichier.
    \item Import des dossiers et des validations de l'A6 (Excel, ODS, CSV).
  \end{itemize}
\end{frame}

% ---------------------------------------------------------------------------
\begin{frame}{Un logiciel libre, testé en continu}
  \begin{columns}[T]
    \begin{column}{0.48\textwidth}
      \begin{itemize}
        \item Moodle 4.0 et plus, PHP 8.
        \item Licence \alert{GPL v3}.
        \item Code au standard Moodle, vérifié à chaque modification.
        \item Guide utilisateur PDF publié avec chaque version.
      \end{itemize}
    \end{column}
    \begin{column}{0.48\textwidth}
      \textbf{Intégration continue}
      \begin{itemize}
        \item \alert{211 tests} PHPUnit et 4 scénarios Behat.
        \item Moodle 4.1 (MariaDB) et 4.5 (PostgreSQL).
        \item Mises à jour de base testées.
        \item Publication automatique des versions.
      \end{itemize}
    \end{column}
  \end{columns}
\end{frame}

% ---------------------------------------------------------------------------
\begin{frame}{Mise en place}
  \begin{enumerate}
    \item Copier \texttt{mod/these} et \texttt{mod/thesesynthesis} dans Moodle, lancer la mise à jour.
    \item Créer un cours par promotion avec une activité \og Gestion des thèses\fg{}, et un cours \og Suivi des thèses\fg{}.
    \item Inscrire étudiants, directeurs (Enseignant non éditeur), DEVE (Enseignant), bibliothèque.
    \item Renseigner les \alert{paramètres du circuit} et choisir les agents.
    \item Remplir l'annuaire du jury (import CSV).
  \end{enumerate}
  \vfill
  \small Serveur de test : \texttt{tools/testserver/seed.php} crée comptes, cours et annuaire à partir d'un fichier.
\end{frame}

% ---------------------------------------------------------------------------
\begin{frame}[plain]
  \centering
  \vfill
  {\Huge\bfseries\color{moodle} ThesisManager}\\[1em]
  {\large Un dossier, un circuit, des documents signés}\\[2em]
  {\small\color{gris} \url{https://github.com/vinlmoe/ThesisManager}}\\[0.5em]
  {\small\color{gris} Guide détaillé : guide PDF publié avec chaque version}
  \vfill
\end{frame}

\end{document}
